% GENERATED FILE. Edit canonical lessons, then run npm run generate:books.
% canonical-source-hash: fnv1a64:8e916d245e578cb8
% canonical-lessons: IT-C527-biografia, IT-C527-fantascienza, IT-C527-saggio, IT-C527-fumetto, IT-C527-poeta

\chapter{Che cosa leggi?}
\label{ch:it-che-cosa-leggi}
\hlchaptermodality{527}

\begin{chapteropening}
\textbf{By the end of this chapter:} \emph{I can talk about the kinds of books I like.}

Complete the last lesson of chapter 527: I can talk about the kinds of books I like.
\end{chapteropening}

\section[la biografia]{la biografia — biography}

\label{lesson:IT-C527-biografia}



\begin{quote}
Before the new one: say the Italian for to devour, then the Italian for to reread.
\end{quote}

\subsection*{You'll want to know: la biografia}
\textbf{la biografia} — \textquotedblleft{}biography\textquotedblright{}.

\begin{quote}
Sto leggendo la biografia di un famoso pittore.
\end{quote}

\begin{quote}
\emph{I am reading the biography of a famous painter.}
\end{quote}

\emph{L'autobiografia} — autobiography.

\begin{cousinweb}
\emph{biografia} is a learned word built from Greek \emph{bíos} \textquotedblleft{}life\textquotedblright{} and \emph{-graphía} \textquotedblleft{}writing\textquotedblright{} (EtymDB).
\end{cousinweb}

\begin{grammarlens}[title={What you've built}]
One more word for telling, explaining and describing at B1.
\end{grammarlens}

\subsection*{Your turn}
\begin{itemize}
\raggedright
  \item \emph{Say it:} \emph{la biografia}
  \item \emph{Say it:} \emph{la biografia}, once more
  \item \emph{Say it:} a sentence of your own with \emph{la biografia}
  \item \emph{Recall:} say the Italian for to devour, then the Italian for to reread, then say \emph{la biografia} again
\end{itemize}

\begin{tcolorbox}[breakable,colback=teal!4,colframe=teal!35!black,title={Before you move on}]
What does la biografia mean? (\textquotedblleft{}Biography\textquotedblright{}.) Say it once more.
\end{tcolorbox}
\section[la fantascienza]{la fantascienza — science fiction}

\label{lesson:IT-C527-fantascienza}



\begin{quote}
Before the new one: say the Italian for to reread, then the Italian for biography.
\end{quote}

\subsection*{You'll want to know: la fantascienza}
\textbf{la fantascienza} — \textquotedblleft{}science fiction\textquotedblright{}.

\begin{quote}
Mio figlio legge solo fantascienza.
\end{quote}

\begin{quote}
\emph{My son only reads science fiction.}
\end{quote}

Built from \emph{fantasia} and \emph{scienza}.

\begin{cousinweb}
\emph{fantascienza} is a modern word built from \emph{fantasia} \textquotedblleft{}fantasy\textquotedblright{}, from Greek \emph{phantasía}, and \emph{scienza} \textquotedblleft{}science\textquotedblright{}, from Latin \emph{scientia} \textquotedblleft{}knowledge\textquotedblright{} (Lewis \& Short).
\end{cousinweb}

\begin{grammarlens}[title={What you've built}]
One more word for telling, explaining and describing at B1.
\end{grammarlens}

\subsection*{Your turn}
\begin{itemize}
\raggedright
  \item \emph{Say it:} \emph{la fantascienza}
  \item \emph{Say it:} \emph{la fantascienza}, once more
  \item \emph{Say it:} a sentence of your own with \emph{la fantascienza}
  \item \emph{Recall:} say the Italian for to reread, then the Italian for biography, then say \emph{la fantascienza} again
\end{itemize}

\begin{tcolorbox}[breakable,colback=teal!4,colframe=teal!35!black,title={Before you move on}]
What does la fantascienza mean? (\textquotedblleft{}Science fiction\textquotedblright{}.) Say it once more.
\end{tcolorbox}
\section[il saggio]{il saggio — essay; non-fiction book}

\label{lesson:IT-C527-saggio}



\begin{quote}
Before the new one: say the Italian for biography, then the Italian for science fiction.
\end{quote}

\subsection*{You'll want to know: il saggio}
\textbf{il saggio} — \textquotedblleft{}essay; non-fiction book\textquotedblright{}.

\begin{quote}
Ha scritto un saggio sulla storia della musica.
\end{quote}

\begin{quote}
\emph{She wrote an essay on the history of music.}
\end{quote}

As an adjective, \emph{saggio} means wise; that word has a different origin.

\begin{cousinweb}
\emph{saggio} \textquotedblleft{}essay\textquotedblright{} comes from Late Latin \emph{exagium} \textquotedblleft{}a weighing, a test\textquotedblright{} (EtymDB): an essay tries out an idea.
\end{cousinweb}

\begin{grammarlens}[title={What you've built}]
One more word for telling, explaining and describing at B1.
\end{grammarlens}

\subsection*{Your turn}
\begin{itemize}
\raggedright
  \item \emph{Say it:} \emph{il saggio}
  \item \emph{Say it:} \emph{il saggio}, once more
  \item \emph{Say it:} a sentence of your own with \emph{il saggio}
  \item \emph{Recall:} say the Italian for biography, then the Italian for science fiction, then say \emph{il saggio} again
\end{itemize}

\begin{tcolorbox}[breakable,colback=teal!4,colframe=teal!35!black,title={Before you move on}]
What does il saggio mean? (\textquotedblleft{}Essay; non-fiction book\textquotedblright{}.) Say it once more.
\end{tcolorbox}
\section[il fumetto]{il fumetto — comic (strip)}

\label{lesson:IT-C527-fumetto}



\begin{quote}
Before the new one: say the Italian for science fiction, then the Italian for essay.
\end{quote}

\subsection*{You'll want to know: il fumetto}
\textbf{il fumetto} — \textquotedblleft{}comic (strip)\textquotedblright{}.

\begin{quote}
Da piccolo leggevo fumetti tutto il giorno.
\end{quote}

\begin{quote}
\emph{As a child I read comics all day.}
\end{quote}

Literally \textquotedblleft{}a little puff of smoke\textquotedblright{}: the speech balloon.

\begin{cousinweb}
\emph{fumetto} is built on \emph{fumo} \textquotedblleft{}smoke\textquotedblright{} (EtymDB), inherited from Latin \emph{fūmus}, with \emph{-etto}: the speech balloon looks like a puff of smoke.
\end{cousinweb}

\begin{grammarlens}[title={What you've built}]
One more word for telling, explaining and describing at B1.
\end{grammarlens}

\subsection*{Your turn}
\begin{itemize}
\raggedright
  \item \emph{Say it:} \emph{il fumetto}
  \item \emph{Say it:} \emph{il fumetto}, once more
  \item \emph{Say it:} a sentence of your own with \emph{il fumetto}
  \item \emph{Recall:} say the Italian for science fiction, then the Italian for essay, then say \emph{il fumetto} again
\end{itemize}

\begin{tcolorbox}[breakable,colback=teal!4,colframe=teal!35!black,title={Before you move on}]
What does il fumetto mean? (\textquotedblleft{}Comic (strip)\textquotedblright{}.) Say it once more.
\end{tcolorbox}
\section[il poeta]{il poeta — poet}

\label{lesson:IT-C527-poeta}



\begin{quote}
Before the new one: say the Italian for essay, then the Italian for comic.
\end{quote}

\subsection*{You'll want to know: il poeta}
\textbf{il poeta} — \textquotedblleft{}poet\textquotedblright{}.

\begin{quote}
Il poeta legge le sue poesie in piazza.
\end{quote}

\begin{quote}
\emph{The poet reads his poems in the square.}
\end{quote}

Masculine but ending in \emph{-a}: \emph{il poeta}, \emph{i poeti}; a woman is \emph{la poetessa}.

\begin{cousinweb}
\emph{poeta} is a learned borrowing from Latin \emph{poēta} (EtymDB), from Greek \emph{poiēt\'{\={e}}s} \textquotedblleft{}a maker\textquotedblright{}, from \emph{poieîn} \textquotedblleft{}to make\textquotedblright{}.
\end{cousinweb}

\begin{grammarlens}[title={What you've built}]
One more word for telling, explaining and describing at B1.
\end{grammarlens}

\subsection*{Your turn}
\begin{itemize}
\raggedright
  \item \emph{Say it:} \emph{il poeta}
  \item \emph{Say it:} \emph{il poeta}, once more
  \item \emph{Say it:} a sentence of your own with \emph{il poeta}
  \item \emph{Recall:} say the Italian for essay, then the Italian for comic, then say \emph{il poeta} again
\end{itemize}

\begin{tcolorbox}[breakable,colback=teal!4,colframe=teal!35!black,title={Before you move on}]
What does il poeta mean? (\textquotedblleft{}Poet\textquotedblright{}.) Say it once more.
\end{tcolorbox}
